% questions.tex —— 2014 年·陕西卷·文科 · 题目内容片段（无答案）
% 由 tools/generate_projects.py 依来源总账自 2014/2014_数学真题_陕西文_试题.md 生成，勿手改（改源 MD 后重跑）


\begin{gknotice}
\textbf{来源说明}：本题面依据固定源 PDF 的两页 A3 题面逐题转录，原卷印刷页码为 913—914。第 913 面为选择题 1—10、填空题 11—15；第 914 面为解答题 16—21。

\textbf{图形说明}：第 4 题程序框图、第 10 题三次函数图象、第 15 题【B】半圆、第 17 题四面体及三视图、第 20 题椭圆与内切圆，均由项目内 TikZ 矢量源重绘；其余各题原页无图。

\textbf{解析归属}：固定源不附答案或解析；配套答案与解析均由本库按本题面独立复算。
\end{gknotice}

\gkpart{一、选择题}

\begin{question}
\gkti{1}已知集合 $M=\{x\mid x\geqslant 0,x\in\mathbf R\}$，$N=\{x\mid x^2<1,x\in\mathbf R\}$，则 $M\cap N=$\mbox{（　　）}

（A）$[0,1]$；（B）$(0,1)$；（C）$(0,1]$；（D）$[0,1)$。

\end{question}
\begin{question}
\gkti{2}函数 $f(x)=\cos\left(2x+\dfrac\pi4\right)$ 的最小正周期是\mbox{（　　）}

（A）$\dfrac\pi2$；（B）$\pi$；（C）$2\pi$；（D）$4\pi$。

\end{question}
\begin{question}
\gkti{3}已知复数 $z=2-\mathrm i$，则 $z\cdot\bar z$ 的值为\mbox{（　　）}

（A）$5$；（B）$\sqrt5$；（C）$3$；（D）$\sqrt3$。

\Needspace{0.4\textheight}
\end{question}
\begin{question}
\gkti{4}根据如图所示的框图，对大于 $2$ 的整数 $N$，输出的数列的通项公式是\mbox{（　　）}
\nopagebreak[4]
\begin{center}% ============================================================
% f2014_shaanxi_arts_q04_flowchart.tex
% 2014 陕西卷（文科）第 4 题：程序框图
% 来源 PDF：普通高考/2014/2014陕西文.pdf
% 原页：第 1 页（A3），印面 913，第 4 题题干下方
% 按原图重绘；判断框内为“$i>N$”，否分支回到“$a_i=2*S$”。
% ============================================================
\begin{tikzpicture}[>=Stealth, line width=0.7pt, font=\small,
  every node/.style={inner sep=2pt},
  term/.style={draw, rounded corners=2pt, minimum width=2.35cm,
    minimum height=0.62cm, align=center},
  proc/.style={draw, minimum width=2.75cm, minimum height=0.62cm, align=center},
  dec/.style={draw, diamond, aspect=1.9, minimum width=2.35cm,
    minimum height=0.95cm, align=center}]
  \node[term] (start) at (0,0) {开始};
  \node[proc] (input) at (0,-0.95) {输入 $N$};
  \node[proc] (init) at (0,-1.9) {$S=1,\ i=1$};
  \node[proc] (body) at (0,-2.85) {$a_i=2*S$};
  \node[proc] (setS) at (0,-3.8) {$S=a_i$};
  \node[proc] (inc) at (0,-4.75) {$i=i+1$};
  \node[dec] (test) at (0,-5.95) {$i>N$};
  \node[proc, minimum width=3.55cm] (out) at (0,-7.1) {输出 $a_1,a_2,\cdots,a_N$};
  \node[term] (end) at (0,-8.05) {结束};

  \draw[->] (start) -- (input);
  \draw[->] (input) -- (init);
  \draw[->] (init) -- (body);
  \draw[->] (body) -- (setS);
  \draw[->] (setS) -- (inc);
  \draw[->] (inc) -- (test);
  \draw[->] (test) -- node[right] {是} (out);
  \draw[->] (out) -- (end);
  % “否”分支：从判断框左侧绕回循环体
  \draw[->] (test.west) -- ++(-1.55,0) |- (body.west);
  \node at (-1.05,-5.72) {否};
\end{tikzpicture}
\end{center}
\nopagebreak[4]
（A）$a_n=2n$；（B）$a_n=2(n-1)$；（C）$a_n=2^n$；（D）$a_n=2^{n-1}$。

\end{question}
\begin{question}
\gkti{5}将边长为 $1$ 的正方形以其一边所在直线为旋转轴旋转一周，所得几何体的侧面积为\mbox{（　　）}

（A）$4\pi$；（B）$3\pi$；（C）$2\pi$；（D）$\pi$。

\end{question}
\begin{question}
\gkti{6}从正方形四个顶点及其中心这 $5$ 个点中，任取 $2$ 个点，则这 $2$ 个点的距离不小于该正方形边长的概率为\mbox{（　　）}

（A）$\dfrac15$；（B）$\dfrac25$；（C）$\dfrac35$；（D）$\dfrac45$。

\end{question}
\begin{question}
\gkti{7}下列函数中，满足“$f(x+y)=f(x)f(y)$”的单调递增函数是\mbox{（　　）}

（A）$f(x)=x^{\frac12}$；（B）$f(x)=x^3$；（C）$f(x)=\left(\dfrac12\right)^x$；（D）$f(x)=3^x$。

\end{question}
\begin{question}
\gkti{8}原命题为“若 $\dfrac{a_n+a_{n+1}}2<a_n$，$n\in\mathbf N_+$，则 $\{a_n\}$ 为递减数列”，关于逆命题、否命题、逆否命题真假性的判断依次如下，正确的是\mbox{（　　）}

（A）真，真，真；（B）假，假，真；（C）真，真，假；（D）假，假，假。

\end{question}
\begin{question}
\gkti{9}某公司 $10$ 位员工的月工资（单位：元）为 $x_1,x_2,\cdots,x_{10}$，其均值和方差分别为 $\bar x$ 和 $s^2$，若从下月起每位员工的月工资增加 $100$ 元，则这 $10$ 位员工下月工资的均值和方差分别为\mbox{（　　）}

（A）$\bar x$，$s^2+100^2$；（B）$\bar x+100$，$s^2+100^2$；（C）$\bar x$，$s^2$；（D）$\bar x+100$，$s^2$。

\Needspace{0.4\textheight}
\end{question}
\begin{question}
\gkti{10}如图，修建一条公路需要一段环湖弯曲路段与两条直道平滑连接（相切），已知环湖弯曲路段为某三次函数图象的一部分，则该函数的解析式为\mbox{（　　）}
\nopagebreak[4]
\begin{center}% ============================================================
% f2014_shaanxi_arts_q10_cubic.tex
% 2014 陕西卷（文科）第 10 题：三次函数弯曲路段与两条直道相切
% 来源 PDF：普通高考/2014/2014陕西文.pdf
% 原页：第 1 页（A3），印面 913，第 10 题题干下方
% 弯曲路段为 y=x^3/2-x^2/2-x（0≤x≤2，虚线），
% 在 O(0,0) 与 y=-x 相切，在 (2,0) 与 y=3x-6 相切；
% 灰色区域为湖面，其下沿为不规则岸线，岸线不是函数图象。
% ============================================================
\begin{tikzpicture}[>=Stealth, line width=0.7pt, font=\small,
  x=1.05cm, y=0.82cm, line cap=round, line join=round]
  % 湖面（灰色区域，下沿为不规则岸线）
  \fill[gray!22]
    (-2.55,3.00) -- (2.95,3.00) -- (2.75,2.70) -- (2.55,2.40)
    -- (2.40,1.85) -- (2.15,1.75) -- (2.00,1.35) -- (1.70,1.25)
    -- (1.45,1.00) -- (1.10,0.95) -- (0.75,0.78) -- (0.40,0.72)
    -- (0.05,0.90) -- (-0.30,1.05) -- (-0.60,1.25) -- (-0.85,1.55)
    -- (-1.10,1.65) -- (-1.30,1.95) -- (-1.55,2.05) -- (-1.75,2.45)
    -- (-1.95,2.55) -- cycle;
  \draw[line width=0.6pt]
    (2.95,3.00) -- (2.75,2.70) -- (2.55,2.40)
    -- (2.40,1.85) -- (2.15,1.75) -- (2.00,1.35) -- (1.70,1.25)
    -- (1.45,1.00) -- (1.10,0.95) -- (0.75,0.78) -- (0.40,0.72)
    -- (0.05,0.90) -- (-0.30,1.05) -- (-0.60,1.25) -- (-0.85,1.55)
    -- (-1.10,1.65) -- (-1.30,1.95) -- (-1.55,2.05) -- (-1.75,2.45)
    -- (-1.95,2.55) -- (-2.55,3.00);
  \node at (0.55,2.05) {湖面};

  % 坐标轴
  \draw[->] (-4.40,0) -- (4.20,0) node[right] {$x$（千米）};
  \draw[->] (0,-2.65) -- (0,3.35) node[above] {$y$（千米）};
  \node[below left] at (0,0) {$O$};

  % 两条直道
  \draw (-3.40,3.40) -- (0,0);
  \draw (2,0) -- (3.30,3.90);
  \node[left] at (-3.15,2.60) {$y=-x$};
  \node[right] at (3.02,3.02) {$y=3x-6$};

  % 弯曲路段（三次函数图象）
  \draw[dashed, thick, domain=0:2, samples=100]
    plot (\x,{0.5*(\x)^3-0.5*(\x)^2-(\x)});

  % 切点与横坐标标注
  \fill (0,0) circle (1.2pt);
  \fill (2,0) circle (1.2pt);
  \node[below right, inner sep=1pt] at (2,0) {$2$};
\end{tikzpicture}
\end{center}
\nopagebreak[4]
（A）$y=\dfrac12x^3-\dfrac12x^2-x$；（B）$y=\dfrac12x^3+\dfrac12x^2-3x$；（C）$y=\dfrac14x^3-x$；（D）$y=\dfrac14x^3+\dfrac12x^2-2x$。

\end{question}
\gkpart{二、填空题}

\begin{question}
\gkti{11}抛物线 $y^2=4x$ 的准线方程为\gkblank{}。

\end{question}
\begin{question}
\gkti{12}已知 $4^a=2$，$\lg x=a$，则 $x=$\gkblank{}。

\end{question}
\begin{question}
\gkti{13}设 $0<\theta<\dfrac\pi2$，向量 $\boldsymbol a=(\sin2\theta,\cos\theta)$，$\boldsymbol b=(1,-\cos\theta)$，若 $\boldsymbol a\cdot\boldsymbol b=0$，则 $\tan\theta=$\gkblank{}。

\end{question}
\begin{question}
\gkti{14}已知 $f(x)=\dfrac{x}{1+x}$，$x\geqslant 0$，若 $f_1(x)=f(x)$，$f_{n+1}(x)=f(f_n(x))$，$n\in\mathbf N_+$，则 $f_{2014}(x)$ 的表达式为\gkblank{}。

\Needspace{0.4\textheight}
\end{question}
\begin{question}
\gkti{15}三选一。

【A】设 $a,b,m,n\in\mathbf R$，且 $a^2+b^2=5$，$ma+nb=5$，则 $\sqrt{m^2+n^2}$ 的最小值为\gkblank{}。

【B】如图，$\triangle ABC$ 中，$BC=6$，以 $BC$ 为直径的半圆分别交 $AB$，$AC$ 于点 $E$，$F$，若 $AC=2AE$，则 $EF=$\gkblank{}。
\nopagebreak[4]
\begin{center}% ============================================================
% f2014_shaanxi_arts_q15b_semicircle.tex
% 2014 陕西卷（文科）第 15 题【B】：以 BC 为直径的半圆
% 来源 PDF：普通高考/2014/2014陕西文.pdf
% 原页：第 1 页（A3），印面 913，第 15 题【B】题干下方
% 图中 BC=6，E、F 在半圆上，AC=2AE；取 b=5、c=(5+√69)/2，
% 使 BC=6 且 ∠A=60°，此时 EF=3。
% ============================================================
\begin{tikzpicture}[>=Stealth, line width=0.7pt, font=\small, x=0.72cm, y=0.72cm]
  \coordinate (B) at (-3,0);
  \coordinate (C) at (3,0);
  \coordinate (O) at (0,0);
  \coordinate (A) at (1.606,4.802);
  \coordinate (E) at (-0.125,2.997);
  \coordinate (F) at (2.533,1.607);

  % 半圆与三角形
  \draw[domain=0:180, samples=120] plot ({3*cos(\x)},{3*sin(\x)});
  \draw (B) -- (A) -- (C) -- cycle;
  \draw (E) -- (F);

  % 点与标签
  \fill (B) circle (1.1pt) node[below left] {$B$};
  \fill (C) circle (1.1pt) node[below right] {$C$};
  \fill (A) circle (1.1pt) node[above] {$A$};
  \fill (E) circle (1.1pt) node[above left] {$E$};
  \fill (F) circle (1.1pt) node[above right] {$F$};
\end{tikzpicture}
\end{center}
\nopagebreak[4]
【C】在极坐标系中，点 $\left(2,\dfrac\pi6\right)$ 到直线 $\rho\sin\left(\theta-\dfrac\pi6\right)=1$ 的距离是\gkblank{}。

\end{question}
\gkpart{三、解答题}

\begin{question}
\gkti{16}$\triangle ABC$ 的内角 $A$，$B$，$C$ 所对的边分别为 $a$，$b$，$c$。

（1） 若 $a$，$b$，$c$ 成等差数列，证明：$\sin A+\sin C=2\sin(A+C)$；

（2） 若 $a$，$b$，$c$ 成等比数列，且 $c=2a$，求 $\cos B$ 的值。

\Needspace{0.4\textheight}
\end{question}
\begin{question}
\gkti{17}四面体 $ABCD$ 及其三视图如图所示，过棱 $AB$ 的中点 $E$ 作平行于 $AD$，$BC$ 的平面分别交四面体的棱 $BD$，$DC$，$CA$ 于点 $F$，$G$，$H$。
\nopagebreak[4]
\begin{center}% ============================================================
% f2014_shaanxi_arts_q17_tetrahedron_views.tex
% 2014 陕西卷（文科）第 17 题：四面体 ABCD 及其三视图
% 来源 PDF：普通高考/2014/2014陕西文.pdf
% 原页：第 1 页（A3），印面 913，第 17 题题干下方
% 由三视图得 DA⊥DB、DA⊥DC、DB⊥DC，且 DA=1、DB=DC=2；
% 原页俯视图只标出纵向尺寸 2，横向尺寸由主视图的 2 得出。
% 左图为立体示意图（平行投影，不按比例）；右图为三视图。
% ============================================================
\begin{tikzpicture}[>=Stealth, line width=0.7pt, font=\small,
  line cap=round, line join=round]
  % ---------- 立体示意图 ----------
  \begin{scope}[scale=1.28]
  \coordinate (A) at (0.10,2.40);
  \coordinate (B) at (-1.50,0.35);
  \coordinate (C) at (2.60,0.75);
  \coordinate (D) at (0.05,0.95);
  \coordinate (E) at (-0.70,1.375); % AB 中点
  \coordinate (F) at (-0.725,0.65); % BD 中点
  \coordinate (G) at (1.325,0.85);  % DC 中点
  \coordinate (H) at (1.35,1.575);  % CA 中点

  % 可见棱
  \draw (A) -- (B) -- (C) -- cycle;
  \draw (E) -- (H);
  % 不可见棱与截面边
  \draw[dashed] (A) -- (D) -- (B);
  \draw[dashed] (D) -- (C);
  \draw[dashed] (E) -- (F) -- (G) -- (H);

  \fill (A) circle (1.1pt) node[above] {$A$};
  \fill (B) circle (1.1pt) node[below left] {$B$};
  \fill (C) circle (1.1pt) node[right] {$C$};
  \fill (D) circle (1.1pt) node[below] {$D$};
  \fill (E) circle (1.1pt) node[left] {$E$};
  \fill (F) circle (1.1pt) node[below left] {$F$};
  \fill (G) circle (1.1pt) node[below] {$G$};
  \fill (H) circle (1.1pt) node[above right] {$H$};
  \end{scope}

  % ---------- 三视图 ----------
  % 主视图：横向 2，纵向 1
  \begin{scope}[shift={(5.05,0.25)}]
    \draw (0,1) -- (0,0) -- (2,0) -- cycle;
    \draw[<->] (-0.42,0) -- (-0.42,1);
    \node[left] at (-0.42,0.5) {$1$};
    \draw[<->] (0,-0.42) -- (2,-0.42);
    \node[below] at (1,-0.42) {$2$};
    \node[below] at (1,-0.85) {主视图};
  \end{scope}

  % 左视图：原页不标尺寸
  \begin{scope}[shift={(7.85,0.25)}]
    \draw (0,1) -- (0,0) -- (2,0) -- cycle;
    \node[below] at (1,-0.45) {左视图};
  \end{scope}

  % 俯视图：纵向 2，横向由主视图给出
  \begin{scope}[shift={(5.05,-4.20)}]
    \draw (0,0) -- (0,2) -- (2,2) -- cycle;
    \draw[<->] (-0.42,0) -- (-0.42,2);
    \node[left] at (-0.42,1) {$2$};
    \node[below] at (1,-0.45) {俯视图};
  \end{scope}
\end{tikzpicture}
\end{center}
\nopagebreak[4]
（1） 求四面体 $ABCD$ 的体积；

（2） 证明：四边形 $EFGH$ 是矩形。

\end{question}
\begin{question}
\gkti{18}在直角坐标系 $xOy$ 中，已知点 $A(1,1)$，$B(2,3)$，$C(3,2)$，点 $P(x,y)$ 在 $\triangle ABC$ 三边围成的区域（含边界）上，且 $\overrightarrow{OP}=m\overrightarrow{AB}+n\overrightarrow{AC}$（$m,n\in\mathbf R$）。

（1） 若 $m=n=\dfrac23$，求 $|\overrightarrow{OP}|$；

（2） 用 $x$，$y$ 表示 $m-n$，并求 $m-n$ 的最大值。

\end{question}
\begin{question}
\gkti{19}某保险公司利用简单随机抽样方法，对投保车辆进行抽样，样本车辆中每辆车的赔付结果统计如下：

\begin{center}
\begin{tabular}{llllll}
\toprule
赔付金额（元） & $0$ & $1000$ & $2000$ & $3000$ & $4000$ \\
\midrule
车辆数（辆） & $500$ & $130$ & $100$ & $150$ & $120$ \\
\bottomrule
\end{tabular}
\end{center}

（1） 若每辆车的投保金额均为 $2800$ 元，估计赔付金额大于投保金额的概率；

（2） 在样本车辆中，车主是新司机的占 $10\%$，在赔付金额为 $4000$ 元的样本车辆中，车主是新司机的占 $20\%$，估计在已投保车辆中，新司机获赔金额为 $4000$ 元的概率。

\Needspace{0.4\textheight}
\end{question}
\begin{question}
\gkti{20}已知椭圆 $\dfrac{x^2}{a^2}+\dfrac{y^2}{b^2}=1$（$a>b>0$）经过点 $\left(0,\sqrt3\right)$，离心率为 $\dfrac12$，左右焦点分别为 $F_1(-c,0)$，$F_2(c,0)$。
\nopagebreak[4]
\begin{center}% ============================================================
% f2014_shaanxi_arts_q20_ellipse_circle.tex
% 2014 陕西卷（文科）第 20 题：椭圆与以 F1F2 为直径的圆
% 来源 PDF：普通高考/2014/2014陕西文.pdf
% 原页：第 2 页（A3），印面 914，第 20 题题干右侧
% 椭圆 x^2/4+y^2/3=1，圆 x^2+y^2=1；
% 直线 l: y=-x/2+√3/3 与椭圆交于 A、B，与圆交于 C、D。
% A(-1.370,1.262)，B(1.947,-0.396)，C(-0.535,0.845)，D(0.997,0.079)。
% ============================================================
\begin{tikzpicture}[>=Stealth, line width=0.7pt, font=\small, x=1.55cm, y=1.55cm]
  % 坐标轴
  \draw[->] (-2.85,0) -- (2.95,0) node[right] {$x$};
  \draw[->] (0,-2.35) -- (0,2.55) node[above] {$y$};
  \node[below left] at (0,0) {$O$};

  % 椭圆与圆
  \draw[domain=0:360, samples=180]
    plot ({2*cos(\x)},{1.732*sin(\x)});
  \draw (0,0) circle (1);
  \fill (-1,0) circle (1.1pt);
  \node[below left] at (-1.03,0) {$F_1$};
  \fill (1,0) circle (1.1pt);
  \node[below right] at (1.03,0) {$F_2$};

  % 直线 l
  \draw (-2.15,1.652) -- (2.55,-0.698);

  % 交点
  \coordinate (A) at (-1.370,1.262);
  \coordinate (B) at (1.947,-0.396);
  \coordinate (C) at (-0.535,0.845);
  \coordinate (D) at (0.997,0.079);
  \fill (A) circle (1.2pt) node[above left] {$A$};
  \fill (B) circle (1.2pt) node[below right] {$B$};
  \fill (C) circle (1.2pt) node[above left] {$C$};
  \fill (D) circle (1.2pt) node[above right] {$D$};
\end{tikzpicture}
\end{center}
\nopagebreak[4]
（1） 求椭圆的方程；

（2） 若直线 $l:y=-\dfrac12x+m$ 与椭圆交于 $A$，$B$ 两点，与以 $F_1F_2$ 为直径的圆交于 $C$，$D$ 两点，且满足 $\dfrac{|AB|}{|CD|}=\dfrac{5\sqrt3}{4}$，求直线 $l$ 的方程。

\end{question}
\begin{question}
\gkti{21}设函数 $f(x)=\ln x+\dfrac mx$，$m\in\mathbf R$。

（1） 当 $m=\mathrm e$（$\mathrm e$ 为自然对数的底数）时，求 $f(x)$ 的极小值；

（2） 讨论函数 $g(x)=f'(x)-\dfrac x3$ 零点的个数；

（3） 若对任意 $b>a>0$，$\dfrac{f(b)-f(a)}{b-a}<1$ 恒成立，求 $m$ 的取值范围。

\end{question}